\section{Judges, State and Private Incentives}
\label{sec:theory}
The economics, political science, and legal literature have thoroughly documented a wide range of extralegal factors that may influence judicial decision-making.\footnote{See \cite{Posner-Judges} for a detailed review in the U.S. context.} These include unobserved differences in judge preferences and characteristics, often captured by judge fixed effects \citep{Kling-Random, Aizer-Juvenile}; electoral cycles and political incentives \citep{Gordon-Judges-Election, Berdejo-PoliticalCycles}; institutional capacity constraints \citep{Schnepel}; and fiscal considerations \citep{Ouss-Incentives}, among others. However, much of this work has focused on the United States and other countries in the Global North, raising concerns about the external validity of these findings for contexts in the Global South, where institutions are often weaker and corruption more prevalent.

In the case of Mexico, as described above, several mechanisms may be driving the observed relationship between private prison construction and sentencing outcomes. Given the specific institutional setting and cost structures, the direction of the estimated effect can offer insight into the dominant mechanism at play. I hypothesize that a positive coefficient—indicating an increase in sentencing following the opening of a nearby private prison—would be consistent with two main explanations.

First, the state may be incentivizing judges to impose longer sentences in order to maximize the use of newly installed capacity. This mechanism is particularly relevant in the Mexican context, where contracts specify payments based on installed capacity rather than per-inmate fees. This is in line with some empirical evidence that suggests that increased capacity by itself leads to higher incarcerations (Parkinson's law) \citep{Parkinson-Stolz}.\footnote{I do not have much on this so any suggestion will be welcomed.}

Second, judges may be responding to capacity constraints in public prisons. In periods of overcrowding, they may have issued more lenient sentences due to limited available space. Once a new facility becomes operational, sentencing may return to prior levels, resulting in an apparent increase. This explanation aligns with the findings of \cite{Ouss-Incentives}, who shows that when prison costs are shifted from the state to county governments, judges—who are accountable at the county level—respond by reducing sentence lengths and incarceration rates. Similarly, \cite{Schnepel} demonstrates that overcrowding in Texas prisons, combined with a failed ballot initiative to expand capacity, led to an increased probability of diversion away from incarceration. \cite{Dippel} find no effect of public prison construction on incarceration or sentence length, but do observe increases when private prisons are built. Although not entirely conclusive nor fully causal, they propose a salience-based mechanism: once new infrastructure becomes visible, communities and officials may expect it to be used, driving up incarceration.

Finally, the possibility that judges are influenced by private actors—whether through lobbying, informal pressure, or outright corruption—cannot be dismissed. In the Mexican context, this channel is particularly interesting given the contractual structure of private prison operations. Since companies are paid based on installed capacity rather than occupancy, their financial incentive may actually favor receiving fewer inmates or inmates with shorter sentences. Anecdotal evidence suggest this could happen.\footnote{The 2011 “Kids for Cash” scandal in Pennsylvania, where two judges were convicted of accepting bribes from private juvenile detention centers in exchange for imposing longer and harsher sentences on minors.\footnote{See: \href{https://jlc.org/luzerne-kids-cash-scandal}{Juvenile Law Center}}} While \cite{Galinato-2020, Dippel} find no consistent evidence that corruption is the primary driver of sentencing increases in their settings, this mechanism remains plausible—especially in countries like Mexico, where judicial systems are more vulnerable to capture and corruption is more widespread.